
\documentclass{article}
\usepackage{colortbl}
\usepackage{makecell}
\usepackage{multirow}
\usepackage{supertabular}

\begin{document}

\newcounter{utterance}

\centering \large Interaction Transcript for game `eqbench', experiment `eq\_bench', episode 124 with reflexion{-}llama{-}3.1{-}8b{-}instruct.
\vspace{24pt}

{ \footnotesize  \setcounter{utterance}{1}
\setlength{\tabcolsep}{0pt}
\begin{supertabular}{c@{$\;$}|p{.15\linewidth}@{}p{.15\linewidth}p{.15\linewidth}p{.15\linewidth}p{.15\linewidth}p{.15\linewidth}}
    \# & \multicolumn{2}{c}{Player} && \multicolumn{2}{c}{Game Master} \\
    \hline

    \theutterance \stepcounter{utterance}  
    & & & \multicolumn{4}{p{0.6\linewidth}}{
        \cellcolor[rgb]{0.9,0.9,0.9}{
            \makecell[{{p{\linewidth}}}]{
                \texttt{\tiny{[P1$\langle$GM]}}
                \texttt{You have attempted this same task before and did not succeed. These are the plans you wrote after those attempts.} \\
\texttt{Trial \#1: New plan: I got stuck in a simplistic analysis of the emotions expressed in the dialogue, focusing on the surface{-}level reactions of Serepta and Argyle. However, I failed to consider the underlying dynamics and power struggles that are driving their interactions. Next time, I will delve deeper into the conversation, examining the subtle cues and underlying motivations that are likely influencing Serepta's emotional state. I will consider factors such as Serepta's tone, body language, and past interactions with Argyle, as well as Argyle's history of disrespect and Serepta's efforts to maintain control. By taking a more nuanced approach, I will be better equipped to accurately predict Serepta's emotional response.} \\
\\ 
\texttt{Trial \#2: New plan: In my previous attempt, I focused on the explicit emotions expressed in the dialogue, but I didn't consider the underlying power dynamics and Argyle's history of disrespect. This time, I will also examine Serepta's emotional response in the context of her past experiences with Argyle and her efforts to maintain control in their shared space. I will consider how Serepta's tone, body language, and past interactions with Argyle might influence her emotional state, and how these factors might interact with the explicit emotions expressed in the dialogue. Specifically, I will look for signs of frustration, resentment, or even anger that might be underlying Serepta's surface{-}level emotions. By taking a more holistic approach, I will be better equipped to accurately predict Serepta's emotional response.} \\
\\ 
\texttt{Trial \#3: New plan: In my previous attempts, I focused on the explicit emotions expressed in the dialogue and the underlying power dynamics between Serepta and Argyle. However, I didn't consider the specific context of their relationship and the impact of Argyle's behavior on Serepta's emotional state. This time, I will take a more detailed approach by examining the specific dynamics of their relationship, including the frequency and severity of Argyle's disrespect, Serepta's efforts to maintain control, and the emotional toll that these interactions have taken on Serepta. I will also consider how Serepta's emotional response might be influenced by her own personality, values, and coping mechanisms. By taking a more nuanced and contextualized approach, I will be better equipped to accurately predict Serepta's emotional response.} \\
\\ 
\\ 
\\ 
\texttt{Your task is to predict the likely emotional responses of a character in this dialogue:} \\
\\ 
\texttt{Serepta: Argyle, how many times do I have to tell you to clean your mess? } \\
\texttt{Argyle: Serepta, it's just a few sketches scattered around. It's not like I'm hoarding garbage.} \\
\texttt{Serepta: It's not about the sketches, Argyle. It's about respect. This is our shared space.} \\
\texttt{Argyle: Respect, huh? Like the respect you show when you rearrange my things without asking?} \\
\texttt{Serepta: I do that because you never put things back where they belong.} \\
\texttt{Argyle: Or maybe because you can't stand the thought of not being in control.} \\
\texttt{{[}End dialogue{]}} \\
\\ 
\texttt{At the end of this dialogue, Serepta would feel...} \\
\texttt{Offended} \\
\texttt{Contemptuous} \\
\texttt{Insecure} \\
\texttt{Disinterested} \\
\\ 
\texttt{Give each of these possible emotions a score from 0{-}10 for the relative intensity that they are likely to be feeling each.} \\
\\ 
\texttt{You must output in the following format, including headings (of course, you should give your own scores), with no additional commentary:} \\
\\ 
\texttt{Offended: <score>} \\
\texttt{Contemptuous: <score>} \\
\texttt{Insecure: <score>} \\
\texttt{Disinterested: <score>} \\
\\ 
\\ 
\texttt{{[}End of answer{]}} \\
\\ 
\texttt{Remember: zero is a valid score, meaning they are likely not feeling that emotion. You must score at least one emotion > 0.} \\
\\ 
\texttt{Your answer:} \\
            }
        }
    }
     \\ \\

    \theutterance \stepcounter{utterance}  
    & \multicolumn{4}{p{0.6\linewidth}}{
        \cellcolor[rgb]{0.561,0.737,0.561}{
            \makecell[{{p{\linewidth}}}]{
                \texttt{\tiny{[P1$\rangle$GM]}}
                \texttt{Offended: 8} \\
\texttt{Contemptuous: 9} \\
\texttt{Insecure: 2} \\
\texttt{Disinterested: 0} \\
            }
        }
    }
    & & \\ \\

    \theutterance \stepcounter{utterance}  
    & & & \multicolumn{2}{p{0.3\linewidth}}{
        \cellcolor[rgb]{0.9,0.9,0.9}{
            \makecell[{{p{\linewidth}}}]{
                \texttt{\tiny{[GM$|$GM]}}
                \texttt{Offended: 8} \\
\texttt{Contemptuous: 9} \\
\texttt{Insecure: 2} \\
\texttt{Disinterested: 0} \\
            }
        }
    }
    & & \\ \\

    \theutterance \stepcounter{utterance}  
    & & & \multicolumn{2}{p{0.3\linewidth}}{
        \cellcolor[rgb]{0.9,0.9,0.9}{
            \makecell[{{p{\linewidth}}}]{
                \texttt{\tiny{[GM$|$GM]}}
                \texttt{Offended: 8} \\
\texttt{Contemptuous: 6} \\
\texttt{Insecure: 0} \\
\texttt{Disinterested: 0} \\
            }
        }
    }
    & & \\ \\

    \theutterance \stepcounter{utterance}  
    & & & \multicolumn{2}{p{0.3\linewidth}}{
        \cellcolor[rgb]{0.9,0.9,0.9}{
            \makecell[{{p{\linewidth}}}]{
                \texttt{\tiny{[GM$|$GM]}}
                \texttt{game\_result = LOSE} \\
            }
        }
    }
    & & \\ \\

    \theutterance \stepcounter{utterance}  
    & & \multicolumn{4}{p{0.6\linewidth}}{
        \cellcolor[rgb]{0.95,0.95,0.95}{
            \makecell[{{p{\linewidth}}}]{
                \texttt{\tiny{[P1$\leftrightarrow$P]}}
                \texttt{Post{-}episode Reflexion (not sent to GM):} \\
\texttt{New plan: In my previous attempts, I focused on the explicit emotions expressed in the dialogue, the underlying power dynamics between Serepta and Argyle, and the specific context of their relationship. However, I didn't consider the potential impact of Serepta's own emotional regulation and coping mechanisms on her emotional response. This time, I will take a more comprehensive approach by examining Serepta's emotional regulation strategies, such as her ability to manage frustration, her tendency to suppress or express emotions, and her coping mechanisms for dealing with stress and conflict. I will also consider how these factors might interact with the explicit emotions expressed in the dialogue, the underlying power dynamics, and the specific context of their relationship. By taking a more holistic and nuanced approach, I will be better equipped to accurately predict Serepta's emotional response.} \\
\\ 
\texttt{Offended: 8} \\
\texttt{Contemptuous: 9} \\
\texttt{Insecure: 2} \\
\texttt{Disinterested: 0} \\
            }
        }
    }
    & \\ \\

\end{supertabular}
}

\end{document}
